
%\documentclass[a4]{amsart}
\documentclass[]{amsart}

\usepackage{amsmath,amsthm,amsfonts,setspace}
\usepackage{harvard,MnSymbol,graphicx}
\usepackage{lscape,longtable}

%\doublespacing

\begin{document}

\title[Dealing with Expections]{Mathematical Expectations: Some
  Information and Practical Problems}

\thanks{I cannot claim to be the author of these notes, as I am taking
this information from a wide range of sources.}  

\author{Steven F. Koch$^\dag$}


%\address{$^\dag$ Professor and Head, Department of Economics,
%  Director, Health and Development Policy Research Group, University
%  of Pretoria, Pretoria, Republic of South Africa; (O) 27-12-420-5285,
%  (F) 27-12-362-5207.} 
%\email{steve.koch@up.ac.za}

\keywords{Expectations, Variances}


%\begin{abstract}
%\end{abstract}

\date{March 2023}
\maketitle

\section{Introduction}
In econometrics, the mathematical expectation and the conditional
expectation are extremely important concepts. For that reason, I am
providing a short discussion about them, as well as putting some
problems together that each of you should be able to do.

\section{The Expectation Operator}
Expectations represent the first moment of a distribution.

\subsection{Discrete Random Variables}
The expectation operator is a mathematical construct developed for
random variables. If we consider a sequence of random variables $x$,
which are contained in the population $X$, then 
\begin{equation}\label{disexp}
E[X] = \sum_{x\in X} xf(x)
\end{equation}
In \eqref{disexp}, $f(x)$ refers to the probability of the specific realization.

Assuming that $a$ is a constant, and that $y$ is another sequence of
random variables contained in $Y$, it should be rather straightforward
to show the following properties.
\begin{itemize}
\item $E[a] = a$
\item $E[aX] = aE[X]$
\item $E[aX + bY] = aE[X] + bE[Y]$
\end{itemize} 

In the case of paired random variables, $X$ and $Y$, it should be
noted that a more general definition is needed.  Below, $g(x,y)$ is
the joint probability of the observation; importantly, the observation
is an ordered pair.
\begin{equation}
E[XY]=\sum_{x\in X} \sum_{y \in Y} xy \cdot g(x,y)
\end{equation}

From the preceding definition, it is possible to show:
\begin{itemize}
\item If $X$ and $Y$ are independent, $E[XY]=E[X]\cdot E[Y]$.
\end{itemize}
This proof is rather complicated, and requires you to make use of
marginal probabilities.

\subsection{Continuous Random Variables}
If variables are continuous, each of the sums in the preceding
definitions can be replaced by the representative integral. 
\begin{equation}\label{contexp}
E[X]=\int_{x \in X} x f(x) dx
\end{equation}
However, rather than referring to $f(x)$ as the probability of a
specific realization, in the continuous case, $f(x)$ refers to the
probability density function.

For pairs of random variables, a similar modification is required.
\begin{equation}
E[XY]=\int_{x\in X} \int_{y \in Y} xy \cdot g(x,y) dydx
\end{equation}

\section{Variances}
Variances reflect the second moment of the distribution of a random
variable. Regardless of whether or not the random variable is
continuous or discrete, the following defintion for the variance is
appropriate. 
\begin{equation}
\label{var}
V(X) = E\left[(X-E[X])^2\right]
\end{equation}

Using the definition from \eqref{var}, you should be able to show the
following results.
\begin{itemize}
\item $V(a) = 0$
\item $V(aX)=a^2V(X)$
\item $V(a+X) = V(X)$
\item $V(X)=E\left[X^2\right]-(E[X])^2$
\end{itemize}

\section{Covariances}
Covariances are also reflected within the second moment of the
distribution. In general, it is possible to use matrix algebra to
develop both the variances and covariances of a complete multivariate
distribution. Intuitively, the covariance definition looks similar to
the variance definition and vice versa.
\begin{equation}\label{cov}
C(X,Y) = E\left[(X-E[X])(Y-E[Y])\right]
\end{equation}

From this definition, a number of properties can be derived.
\begin{itemize}
\item $C(a,X) = 0$
\item $C(X,X) = V(X)$
\item $C(X,Y)=E[XY]-E[X]\cdot E[Y]$
\end{itemize}

Furthermore, results from above, related to expectations, can be used to
show the following results.
\begin{itemize}
\item If $X$ and $Y$ are independent, $C(X,Y)=0$.
\item If $E[\epsilon] = 0$ and $C(\epsilon,X) = 0$, $E[X\epsilon]=0$,
  where $\epsilon$ is a random variable.
\end{itemize}



\end{document}